%% Singlify AI — Development Roadmap
%% Migrated from singlet-intelligence/docs/roadmap.md
\documentclass[10pt,twocolumn,letterpaper]{article}
\usepackage{../../templates/singletai-preprint}
\usepackage{graphicx}
\usepackage{amsmath,amssymb}
\usepackage{booktabs}
\usepackage{natbib}
\bibliographystyle{unsrtnat}

% ── Metadata ─────────────────────────────────────────────────────
\title{Singlify AI Development Roadmap: From Transcriptomics \\
  to Clinical Intelligence}
\author{
  Zach DeBruine\textsuperscript{1,*}
  \\[6pt]
  \textsuperscript{1}Department of Computing, Grand Valley State University, Allendale, MI 49401
}

\begin{document}
\maketitle
\correspondingauthor{debruinz@gvsu.edu}

\begin{abstract}
This document defines the four-phase product roadmap for Singlify AI, a
genomics intelligence platform. Phase~1 (v1.0, Transcriptomics Intelligence)
delivers a generative model over ${\sim}100$M human single-cell transcriptomes
via NMF-based gene programs and a Conditional Program Model (CPM) encoder.
Phase~2 (v2.0) expands to ${\sim}1$B cells with multi-modal prediction
(ATAC, CITE-seq, spatial) and RNA velocity dynamics. Phase~3 (v3.0, Genomics
Intelligence) bridges DNA sequence to single-cell expression via
sequence-to-function models trained on GTEx and TenK10K data. Phase~4 (v4.0,
Clinical Intelligence) enables variant-to-clinical-consequence reasoning with
Epic EHR integration. Each phase includes quantitative success metrics,
validation protocols, and deployment milestones.
\end{abstract}

% ── Phase 1 ──────────────────────────────────────────────────────
\section{Phase 1: v1.0 --- Transcriptomics Intelligence}

\subsection{Data Foundation: 100M Human Transcriptomes}

All data comes from NCBI GEO, reprocessed from raw FASTQ through the
\texttt{scgeo} pipeline. This produces:
\begin{itemize}
  \item Spliced/Unspliced/Ambiguous (S/U/A) count layers via simpleaf/alevin-fry
  \item Kraken2 non-host classification per cell (bacteria, virus, fungi)
  \item Unified gene references (GENCODE v44 human)
  \item Consistent QC (doublet scoring, ambient RNA, mitochondrial fraction)
\end{itemize}

The initial 100M cells are human-only, droplet-based (10x Chromium, Drop-seq,
inDrop). Plate-based protocols are excluded from v1.0.

\paragraph{Metadata annotations.} Three sources: (1)~automated cell typing via
SingleR against curated atlases, (2)~LLM extraction from GEO series
descriptions standardized to controlled vocabularies (Uberon, Cell Ontology,
MONDO, HsapDv, ChEMBL/DrugBank), and (3)~pipeline QC metrics.

\subsection{NMF Model Training}

\paragraph{Rank determination.} Speckled holdout cross-validation: mask
${\sim}1\%$ of nonzero entries, train NMF for each candidate rank
$k \in \{256, 512, 1024, 2048, 4096\}$, select the elbow. Ranks constrained
to powers of~2 for GPU efficiency.

\paragraph{Hierarchical depth.} Evaluate 1--3 layers of hierarchical NMF
($X \approx W_1 H_1$ vs.\ $X \approx W_1 W_2 H_2$ vs.\ three layers).
Selection by held-out reconstruction at matched parameter count.

\paragraph{Regularization.} Two hyperparameters:
\begin{equation}
  \min_{W,H \geq 0} \|X - WH\|_F^2
  + \lambda_{\text{L1}}(\|W\|_1 + \|H\|_1)
  + \lambda_{\text{ang}} \sum_{j \neq j'} \max(0, \cos(W_{:,j}, W_{:,j'}) - \epsilon)^2
\end{equation}
Angular regularization prevents degenerate duplicate programs at high rank.
L1 promotes sparse, interpretable programs. Selection: grid search on 10\%
subsample, Pareto-optimal on reconstruction vs.\ revealable rank vs.\ pathway
enrichment specificity.

\paragraph{Final training.} Spliced count matrix
$X_S \in \mathbb{R}_{\geq 0}^{G \times N}$ ($G \approx 30{,}000$,
$N \approx 10^8$) in StreamPress \texttt{.spz} format. CUDA streaming
out-of-core multiplicative updates. Convergence: relative change $< 10^{-6}$
for 5 iterations, or max 500 iterations. Estimated compute: 8$\times$ H100
GPUs, 30--130 GPU-hours.

\subsection{Factor Annotation}

For each of $k$ programs: (1)~enrichment against GO, Reactome, KEGG, MSigDB
Hallmark; (2)~cell type, tissue, disease, and covariate enrichment;
(3)~LLM-generated names and descriptions. Quality target: ${>}80\%$ have
meaningful non-generic names.

\subsection{Conditional Program Encoder (CPM)}

\paragraph{Architecture.} Every metadata field embedded independently
(512-dim for cell type/tissue/disease/perturbation, 128-dim for sex/assay,
256-dim for donor). All embeddings attend via multi-head cross-attention, then
MLP backbone to 256-dim bottleneck $z$. Three output heads: gate ($\alpha$),
magnitude ($\mu$), spread ($\sigma$).

\paragraph{Spike-and-slab reparameterization.} $h_j = g_j \cdot m_j$ where
$g_j \in \{0,1\}$ (hard concrete gate) and $m_j = \exp(\mu_j + \sigma_j
\epsilon_j)$ (log-normal magnitude). Produces exact zeros for inactive
programs, matching NMF sparsity (Gini ${>}0.7$).

\paragraph{Loss function.}
\begin{equation}
  \mathcal{L} = \sum_{j=1}^{k} \ell_j
  + \lambda_{\text{sparse}} \sum_j \pi_j
  + \lambda_{\text{consist}} \mathcal{L}_{\text{consist}}
  + \lambda_{\text{diversity}} \mathcal{L}_{\text{diversity}}
\end{equation}
where $\ell_j$ is spike-and-slab NLL (gate + magnitude for active programs,
gate-only for inactive).

\paragraph{Validation.} Structured holdout at three levels: interpolation
(random cells within known combos), single-axis extrapolation (hold out one
metadata value), and combinatorial extrapolation (hold out multi-axis combos).
Targets: cosine similarity ${>}0.8$, gate accuracy ${>}85\%$, combinatorial
degradation ${<}15\%$.

\subsection{Prediction Confidence Scores}

Every prediction carries empirical confidence grounded in cross-validation
performance. During training, systematic holdouts produce an accuracy
distribution per metadata distance. At inference:
\begin{enumerate}
  \item Compute nearest training neighborhood (embedding distance)
  \item Retrieve empirical accuracy at that distance
  \item Convert to p-value against null distribution
\end{enumerate}

Calibration requirement: 90\% CI covers 90\% $\pm$ 3\% of observations
(Platt scaling + temperature scaling). This is a v1.0-beta launch requirement.

\subsection{Generative Capabilities}

v1.0 supports: (1)~expression profile prediction for fully specified
conditions, (2)~per-gene expression across cell types, (3)~perturbation
effect prediction via $\Delta h$, (4)~disease extrapolation across tissues,
(5)~synthetic reference panel generation.

\subsection{AI Reasoning Layer}

Natural language $\to$ structured CPM queries via LLM parsing. MCP tool
exposure enables AI agents to call \texttt{predict\_expression} programmatically.
Plot code generation returns executable Python/R visualizations alongside
AI interpretation.

\subsection{v1.0 Success Metrics}

\begin{table}[h]
\centering\small
\begin{tabular}{@{}lr@{}}
\toprule
Metric & Target \\
\midrule
Speckled holdout $R^2$ & ${>}0.3$ \\
Meaningful program annotations & ${>}80\%$ of $k$ \\
Encoder combinatorial generalization & ${<}15\%$ degradation \\
Gate accuracy & ${>}85\%$ \\
Confidence calibration & 90\% CI $\pm$ 3\% \\
Cross-species transfer (v1.3+) & ${>}70\%$ accuracy \\
Query latency (point) & ${<}200$~ms \\
Query latency (100 cells) & ${<}2$~s \\
\bottomrule
\end{tabular}
\end{table}

% ── Phase 2 ──────────────────────────────────────────────────────
\section{Phase 2: v2.0 --- Expanded Transcriptomics}

\subsection{Expanded Data Corpus}

Scale to ${\sim}500$M--1B human cells including plate-based protocols
(Smart-seq2, CEL-seq2, MARS-seq). Assay-aware NMF training handles
distributional shift. Enhanced metadata: ancestry/population, detailed
perturbation parsing, deeper developmental staging, author-annotated cell
types.

\subsection{Multi-Modal Prediction}

\paragraph{ATAC-seq.} Reprocess paired scRNA-seq + scATAC-seq (10x Multiome),
train NMF on peak matrices, learn cross-modal projection from RNA to ATAC
program space. Deliverable: predict chromatin accessibility for any condition.

\paragraph{CITE-seq.} Predict surface protein expression from condition
metadata via RNA $\to$ protein program mapping. Relevant to flow cytometry
panel design and CAR-T target selection.

\paragraph{Spatial (Visium).} Deconvolve spots into cell type proportions via
NMF projection. Learn spatial neighbor prediction: given a cell's programs,
predict adjacent cells' programs---generating tissue neighborhoods.

\subsection{Splicing Dynamics and Next-Cell Prediction}

Factor velocities: $V_j = H_{U,j} - \gamma \odot H_{S,j}$. Next-cell
prediction: train the CPM encoder with a \texttt{next\_state} output head to
chain trajectory predictions $h_0 \to h_1 \to \cdots \to h_T$.

\subsection{Composition Head}

$\pi = \text{softmax}(W_\pi \cdot z)$ predicts cell type proportions for
underspecified queries. Label dropout ($p \sim 0.1$--$0.3$) during training
teaches prediction from partial metadata.

\subsection{v2.0 Success Metrics}

\begin{table}[h]
\centering\small
\begin{tabular}{@{}lr@{}}
\toprule
Metric & Target \\
\midrule
RNA $\to$ ATAC (Pearson $r$ per peak) & ${>}0.6$ \\
Spatial neighbor (top-5 accuracy) & ${>}40\%$ \\
Next-cell cosine similarity & ${>}0.7$ \\
Composition KL divergence & ${<}0.1$ \\
Multi-modal AI query success & ${>}80\%$ \\
\bottomrule
\end{tabular}
\end{table}

% ── Phase 3 ──────────────────────────────────────────────────────
\section{Phase 3: v3.0 --- Genomics Intelligence}

\subsection{Sequence-to-Function Foundation}

Train a convolutional + transformer model on GTEx (${\sim}1{,}000$ donors,
54 tissues) + ENCODE/Roadmap to learn regulatory grammar from DNA
($\pm500$~kb around each gene) $\to$ expression/accessibility/splicing.
Alternative: fine-tune Borzoi (Apache 2.0, open weights).

\subsection{Single-Cell Prediction Head}

TenK10K dataset (1,925 donors, paired WGS + scRNA-seq, ${\sim}5$M+ cells).
Architecture: sequence trunk $\to$ condition on cell type + tissue embeddings
(from CPM vocabulary) $\to$ predict program activities $h$ in NMF space.

\textbf{Novel capability:} For any genetic variant, predict cell-type-specific
transcriptional perturbation: not ``this gene changes expression'' but
``Program~847 (hepatocyte cholesterol biosynthesis) is disrupted in
hepatocytes.''

\subsection{v3.0 Success Metrics}

\begin{table}[h]
\centering\small
\begin{tabular}{@{}lr@{}}
\toprule
Metric & Target \\
\midrule
GTEx genotype $\to$ expression (Pearson $r$) & ${>}0.7$ \\
TenK10K scRNA prediction (cosine sim) & ${>}0.6$ \\
ClinVar variant mechanism match & ${>}60\%$ \\
Rare disease blood $R^2$ & ${>}0.4$ \\
\bottomrule
\end{tabular}
\end{table}

% ── Phase 4 ──────────────────────────────────────────────────────
\section{Phase 4: v4.0 --- Clinical Intelligence}

\subsection{Clinical Variant Reasoning}

A clinical LLM with v1.0--v3.0 tools as MCP endpoints. Given a novel variant:
(1)~genomic prediction of cell-type-specific expression changes,
(2)~program disruption analysis, (3)~literature grounding, (4)~clinical
reasoning synthesis.

Training on published clinical case reports: extract variant, run prediction
stack, compare to actual case presentation. Validation on ${\sim}1{,}000$
held-out PubMed cases scored by expert clinicians.

\subsection{Epic EHR Integration}

Epic App Orchard application: VUS triggers analysis $\to$ mechanistic
explanation embedded in patient chart. Initial launch as ``information only''
(Clinical Decision Support exemption under 21st Century Cures Act).

\subsection{v4.0 Success Metrics}

\begin{table}[h]
\centering\small
\begin{tabular}{@{}lr@{}}
\toprule
Metric & Target \\
\midrule
Clinical detail prediction accuracy & ${>}70\%$ \\
VUS reclassification rate & ${>}20\%$ \\
Clinician NPS & ${>}50$ \\
Time to diagnosis (rare disease) & ${<}1$ year \\
\bottomrule
\end{tabular}
\end{table}

% ── Timeline ─────────────────────────────────────────────────────
\section{Execution Timeline}

\begin{table}[h]
\centering\small
\begin{tabular}{@{}lll@{}}
\toprule
Date & Version & Milestone \\
\midrule
2026 Q2 & v1.0-$\alpha$ & NMF model + factor annotations \\
2026 Q3 & v1.0-$\beta$ & CPM encoder + confidence scores \\
2026 Q4 & v1.0 & AI reasoning layer + MCP tools \\
2027 Q1 & v1.1 & Composition head \\
2027 Q1--Q2 & v1.2 & Expanded data (${\sim}300$M cells) \\
2027 Q2--Q3 & v1.3--v1.4 & Cross-species alignment (201 species) \\
2027 Q3--Q4 & v2.0 & Multi-modal + splicing dynamics \\
2027 Q3--Q4 & v3.0 & Sequence-to-function bridge \\
2028 Q1+ & v4.0 & Clinical intelligence + Epic \\
\bottomrule
\end{tabular}
\end{table}

% ── Risk Factors ─────────────────────────────────────────────────
\section{Risk Factors}

\paragraph{v1.0.} NMF rank may be suboptimal at 100M cells (re-run at v1.2).
Encoder overfitting mitigated via holdout testing + label noise. SingleR
annotations noisy (cross-reference with LLM-extracted labels). AI layer
hallucinations mitigated via grounded reasoning (every claim traced to a
specific program + enrichment). Cross-species adapter may underfit divergent
programs (start frozen human, escalate to joint training).

\paragraph{v2.0.} Plate-based integration risk (assay-aware encoder terms).
Cross-modal prediction noise (start with high-quality paired 10x Multiome).

\paragraph{v3.0.} TenK10K is blood-dominated (use GTEx for tissue priors).
Sequence model requires massive compute (use pre-trained Borzoi trunk).

\paragraph{v4.0.} Regulatory uncertainty (launch as information-only CDS).
Clinical validation is expensive (academic partnerships for shared cost).

\bibliography{../../shared/refs}
\end{document}
